% Divide-and-conquer half of the report.
% Preamble needs: amsmath, amssymb, algorithmic, graphicx, cite, and
%   \newtheorem{lemma}{Lemma}  \newtheorem{theorem}{Theorem}
% \dcdir is the folder holding figures/ (empty when compiled from this folder).

\section{Divide and Conquer: Capped Map Search}
\label{sec:dc}

\subsection{Problem}
\label{sec:dc-problem}

Airbnb's map search returns the listings inside the visible map window, but
never more than a fixed number of them: a crowded window is silently
truncated. A market study of a city needs \emph{every} listing, so a scraper
must choose map windows small enough that none is truncated, using as few
searches as possible. This problem came up while scraping Miami, FL and
Rinc\'on, PR~\cite{danial2023airbnb}; a single city-wide window returned only a
small fraction of the listings.

\subsection{Abstraction}
\label{sec:dc-abstraction}

Let $P$ be a set of $n$ \emph{listings}, each with a location in
$\mathbb{R}^2$. Several listings may share a location, since Airbnb rounds
coordinates; let $m$ be the largest number sharing one. A \emph{window} is a
half-open rectangle $r=[x_0,x_1)\times[y_0,y_1)$, and $P_r$ denotes the
listings located in $r$. The search region $R$ is a window with $P_R=P$.

The search service is an oracle with cap $k\ge1$:
\begin{equation}
Q_k(r)=\begin{cases}
P_r & \text{if } |P_r|\le k,\\
\textsc{overflow} & \text{otherwise.}
\end{cases}
\label{eq:oracle}
\end{equation}
The problem is to compute $P$ from $R$ and $Q_k$ using as few oracle calls
(web requests) as possible.

Splitting a window at its midpoint, both horizontally and vertically, gives
four equal \emph{quadrants}. Splitting recursively yields a
\emph{quadtree}~\cite{finkel1974quad,samet1984quadtree} whose leaves are the
windows actually read.

\subsection{Algorithm}
\label{sec:dc-algorithm}

\textsc{Collect} (Fig.~\ref{fig:dc-alg}) queries a window. If it is not
truncated, its listings are the answer. Otherwise the window is divided into
four quadrants, each is solved recursively, and the answers are united. The
full answer is $\textsc{Collect}(R)$.

\begin{figure}[t]
\begin{algorithmic}[1]
\STATE \textbf{function} \textsc{Collect}$(r)$
\STATE $A\gets Q_k(r)$
\IF{$A\neq\textsc{overflow}$}
  \RETURN $A$ \COMMENT{window fits under the cap}
\ENDIF
\STATE $r_1,\dots,r_4\gets$ quadrants of $r$ \COMMENT{divide}
\RETURN $\bigcup_{i=1}^{4}\textsc{Collect}(r_i)$ \COMMENT{conquer, combine}
\end{algorithmic}
\caption{Enumerating the listings in window $r$ through a capped oracle $Q_k$.}
\label{fig:dc-alg}
\end{figure}

% Placed here in the source so LaTeX can set it at the top of the next page;
% a figure* is always deferred to the page after the one where it appears.
\begin{figure*}[!t]
\centering
\includegraphics[width=\textwidth]{\dcdir figures/scaling}
\caption{Scaling with $n$ at $k=50$, for uniform and clustered listings.
(a)~Running time with the least-squares fit $c\,n\log_2(n/k)$
($R^2=0.995$ uniform, $0.996$ clustered). (b)~Running time divided by
$n\log_2(n/k)$; a flat curve confirms the growth rate. (c)~Requests issued,
against the lower bound $n/k$.}
\label{fig:dc-scaling}
\end{figure*}

\subsection{Correctness}
\label{sec:dc-correctness}

\begin{lemma}\label{lem:partition}
The four quadrants of a window $r$ partition $r$.
\end{lemma}
\begin{IEEEproof}
Let $(x_m,y_m)$ be the midpoint of $r$. For $(x,y)\in r$, exactly one of
$x<x_m$ and $x\ge x_m$ holds, and exactly one of $y<y_m$ and $y\ge y_m$; the
pair picks exactly one quadrant. Each quadrant lies in $r$ since
$x_0\le x_m\le x_1$ and $y_0\le y_m\le y_1$.
\end{IEEEproof}

Half-open windows are what make this work: a listing on a split line falls in
exactly one quadrant, so it is neither lost nor counted twice.

Scale $R$ to the unit square and let $\delta\in(0,1]$ be the smallest
coordinate-wise distance, $\max(|\Delta x|,|\Delta y|)$, between two listings at
different locations ($\delta=1$ if there is only one location).

\begin{lemma}\label{lem:height}
If $m\le k$, the recursion tree of $\textsc{Collect}(R)$ has height at most
$\lceil\log_2(1/\delta)\rceil$; in particular the algorithm terminates.
\end{lemma}
\begin{IEEEproof}
A window at depth $d$ has side $2^{-d}$, so two listings inside it at different
locations are less than $2^{-d}$ apart, i.e.\ $\delta<2^{-d}$. A window is split
only if it holds more than $k\ge m$ listings, hence listings at two or more
locations, hence $d<\log_2(1/\delta)$. So every split window has depth at most
$\lceil\log_2(1/\delta)\rceil-1$ and every window depth at most
$\lceil\log_2(1/\delta)\rceil$. A tree of bounded height with four children per
internal node is finite.
\end{IEEEproof}

\begin{theorem}\label{thm:correct}
If $m\le k$, $\textsc{Collect}(r)=P_r$ for every window $r$, and the leaf
windows partition $r$ with at most $k$ listings each.
\end{theorem}
\begin{IEEEproof}
Strong induction on the height $h$ of the recursion tree at $r$ (finite by
Lemma~\ref{lem:height}). If $h=0$, the oracle did not overflow, so
$|P_r|\le k$ and line~4 returns $Q_k(r)=P_r$; the only leaf is $r$. If $h>0$,
the algorithm recurses on the quadrants $r_1,\dots,r_4$, whose subtrees are
shorter than $h$, so by induction $\textsc{Collect}(r_i)=P_{r_i}$ and the leaves
under $r_i$ partition $r_i$ with at most $k$ listings each. By
Lemma~\ref{lem:partition} each listing of $P_r$ lies in exactly one $r_i$, so
$\bigcup_i P_{r_i}=P_r$ and the leaves of the four subtrees partition $r$.
\end{IEEEproof}

The condition $m\le k$ is also necessary: if more than $k$ listings share a
location, every window containing it overflows, so no strategy using $Q_k$
can reveal them.

Lemmas~\ref{lem:partition}--\ref{lem:height} and Theorem~\ref{thm:correct}
are also machine-checked in Lean~4~\cite{lean4} (Appendix~\ref{app:dc-lean}).
Reported coordinates have finite precision (Airbnb gives five decimals), so
the formalization works on an integer grid, where $\delta$ is one grid unit.
There Lemma~\ref{lem:height} reads: if $m\le k$, recursion depth
$\lceil\log_2 W\rceil$ suffices for a window of side $W$; combined with
Theorem~\ref{thm:correct} this gives total correctness. The real-valued form
of Lemma~\ref{lem:height} is proved only on paper. The nearest existing
formalizations, region quadtrees in Isabelle~\cite{nipkow2024quad,nipkow2024afp}
and $k$-d trees in Isabelle and Coq~\cite{rau2019kdtree,hamid2023kdtree},
obtain termination from a fixed grid exponent or from the point count
shrinking at each median split; neither covers a space-splitting tree over
arbitrary points, so the proof here was written from scratch.

\subsection{Running Time}
\label{sec:dc-analysis}

Let $N$ be the number of oracle calls and $T$ the total work when a node with
$|P_r|$ listings spends $\Theta(|P_r|+1)$ time returning or distributing them.
Let $h$ be the tree height.

Windows at one depth are disjoint (Lemma~\ref{lem:partition}), so each level
holds at most $n$ listings and at most $\lfloor n/(k+1)\rfloor$ split windows,
each holding at least $k+1$ listings. Internal nodes sit at depths
$0,\dots,h-1$ and have four children each, so
\begin{equation}
\lceil n/k\rceil\le N\le 1+4h\Big\lfloor\frac{n}{k+1}\Big\rfloor,
\qquad T=O\big(n(h+1)\big).
\label{eq:general}
\end{equation}
The lower bound holds because the leaves cover $P$ with at most $k$ listings
each. By Lemma~\ref{lem:height}, $T=O(n\log(1/\delta))$: the cost is governed
by how tightly listings crowd, not by $n$ alone.

In the balanced case, where each split sends about $n/4$ listings to each
quadrant (true in expectation for uniformly random locations),
\begin{equation}
T(n)=4\,T(n/4)+\Theta(n)\quad(n>k),
\label{eq:recurrence}
\end{equation}
with $T(n)=\Theta(1)$ for $n\le k$. With $a=b=4$ and
$f(n)=\Theta(n)=\Theta(n^{\log_b a})$, case~2 of the Master
Theorem~\cite{cormen2022introduction} gives $T(n)=\Theta(n\log n)$. Unrolling
the recursion shows the dependence on $k$: the tree is complete with
$L=\lceil\log_4(n/k)\rceil$ levels below the root, each costing $\Theta(n)$, so
\begin{align}
T(n)&=\Theta\!\left(n\log\frac{n}{k}\right),\label{eq:balanced}\\
N&=\sum_{i=0}^{L}4^i=\frac{4^{L+1}-1}{3}=\Theta\!\left(\frac{n}{k}\right),
\label{eq:requests}
\end{align}
within a constant factor of the $n/k$ requests any strategy must make.

\subsection{In the Language of the Problem}
\label{sec:dc-domain}

Search the whole city in one window. If the results fit on the pages the site
shows, record them. Otherwise cut the window into four equal tiles and repeat
on each. A listing exactly on a cut goes to the tile right of or above it, so
it is recorded once. Dense neighborhoods get zoomed into repeatedly; open water
is dismissed after one search. The number of searches is proportional to
$n/k$, the number of listings divided by the number shown per search, by
\eqref{eq:requests}. The only obstacle is an address with more listings than
the cap, such as an apartment building whose units share rounded coordinates:
no zoom level separates them, which is why Theorem~\ref{thm:correct} needs
$m\le k$. In Rinc\'on at most $7$ listings share a location, so any $k\ge7$
works.

\subsection{Experiments}
\label{sec:dc-experiments}

\begin{figure}[!t]
\centering
\includegraphics[width=\columnwidth]{\dcdir figures/depth_vs_delta}
\caption{Effect of packing $k+1$ listings into a square of side $\delta$
($n=10{,}000$, $k=50$): (a)~tree height, (b)~requests.}
\label{fig:dc-depth}
\end{figure}

\begin{figure}[!t]
\centering
\includegraphics[width=2.45in]{\dcdir figures/rincon_quadtree}
\caption{Windows chosen for 261 listings in Rinc\'on, PR ($k=20$).}
\label{fig:dc-rincon}
\end{figure}

\textbf{Setup.} \textsc{Collect} was implemented in Python~3.14
(Appendix~\ref{app:dc-code}) and run on an AMD Ryzen~7 9800X3D. The oracle is
simulated: each node partitions the listings it holds among its quadrants,
which is the $\Theta(|P_r|)$ step of \eqref{eq:recurrence}. We used $k=50$
and $n=2^{10},\dots,2^{19}$ listings in the unit square, either uniform or
drawn from $20$ Gaussian clusters ($\sigma=0.02$). Each size was run five
times on fresh random inputs; we report the fastest run with garbage
collection disabled, as \texttt{timeit} does~\cite{python-timeit}. Every run
was checked to return each listing exactly once.

\textbf{Running time.} Fig.~\ref{fig:dc-scaling}(a) shows measured time
against the fitted model; Fig.~\ref{fig:dc-scaling}(b) divides time by
$n\log_2(n/k)$. Over a $512\times$ range of $n$ the ratio stays within
$26$--$43$\,ns (uniform) and $36$--$48$\,ns (clustered), confirming
$\Theta(n\log(n/k))$. The zig-zag arises because the tree gains a level only
when $n$ grows about fourfold, so the true cost rises in steps; the slight rise
at the largest sizes is consistent with cache effects. Clustered inputs cost
about $20\%$ more because their trees are two levels taller.

\textbf{Requests.} Table~\ref{tab:dc-requests} compares requests on uniform
inputs with the exact balanced count \eqref{eq:requests}; they agree within
$3.1\%$. Fig.~\ref{fig:dc-scaling}(c) shows the count growing linearly,
parallel to the lower bound $n/k$.

\begin{table}[t]
\caption{Requests on uniform inputs ($k=50$) vs.\ the prediction \eqref{eq:requests}}
\label{tab:dc-requests}
\centering
\begin{tabular}{rrrr}
\hline
$n$ & measured & predicted & height \\
\hline
1{,}024   & 85     & 85     & 3 \\
4{,}096   & 333    & 341    & 4 \\
16{,}384  & 1{,}337  & 1{,}365  & 5 \\
65{,}536  & 5{,}293  & 5{,}461  & 6 \\
262{,}144 & 21{,}197 & 21{,}845 & 7 \\
524{,}288 & 21{,}917 & 21{,}845 & 8 \\
\hline
\end{tabular}
\end{table}

\textbf{Crowding.} To test \eqref{eq:general}, we fixed $n=10{,}000$ uniform
listings and added $k+1=51$ listings inside a square of side
$\delta=2^{-2},\dots,2^{-28}$. Fig.~\ref{fig:dc-depth} shows the tree height
tracking $\log_2(1/\delta)$ (Lemma~\ref{lem:height}) although $n$ is fixed,
and each extra level adding about four requests, the $4h$ term of
\eqref{eq:general}.

\textbf{Real data.} Fig.~\ref{fig:dc-rincon} runs the algorithm with $k=20$ on
$261$ Airbnb listings from Rinc\'on, PR (scraped March 29, 2023~\cite{danial2023airbnb}).
It used $57$ requests, $43$ windows and height $5$, within the bounds
$14\le N\le 241$ of \eqref{eq:general}. The two dense clusters are split
repeatedly; the outskirts are covered by a few large windows.
